\section{Additional results on positive maps and channels}

\begin{definition}[Associated positive linear map]
    \label{def:positive_map_is_positive_linear_map}
    \lean{IsPositiveMap.toPositiveLinearMap}
    \leanok
    \uses{def:positive_map}
    A positive matrix map $E : \MN{n} \to \MN{m}$ determines the same
    linear map regarded as a positive linear map: if
    $A,B\in\MN{n}$ and $A \le B$, then $E(A) \le E(B)$ in $\MN{m}$.
\end{definition}

\begin{definition}[Trace-preserving map]\label{def:trace_preserving}
    \lean{IsTracePreservingMap}
    \leanok
    A linear map $E : \MN{D} \to \MN{D}$ is \emph{trace-preserving}
    if $\tr(E(X)) = \tr(X)$ for all $X$.
\end{definition}

\begin{theorem}[Transposition is positive and trace-preserving]
    \label{thm:transpose_positive_trace_preserving}
    \lean{Matrix.transposeLinearMapComplex_isPositiveMap,
        Matrix.transposeLinearMapComplex_isTracePreservingMap}
    \leanok
    \uses{def:positive_map, def:trace_preserving}
    The matrix transposition map $\theta(X)=X^T$ on $\MN{D}$ is positive
    and trace-preserving.
\end{theorem}

\begin{proof}\leanok
    If $X \ge 0$, write $X=C^\dagger C$. Then
    $X^T=C^T\overline C=\overline C^\dagger\overline C$, so $X^T \ge 0$.
    Trace preservation is the identity $\tr(X^T)=\tr(X)$.
\end{proof}

\begin{definition}[Reduction map]
    \label{def:reduction_map}
    \lean{Matrix.reductionMap}
    \leanok
    The reduction map on $\MN{D}$, for $k\in\N$, is
    \[
        T_k(X)=\tr(X)\Id-k^{-1}X.
    \]
\end{definition}

\begin{theorem}[The first reduction map is positive]
    \label{thm:reduction_map_one_positive}
    \lean{Matrix.reductionMap_one_isPositiveMap}
    \leanok
    \uses{def:positive_map, def:reduction_map}
    For $D\ge 1$, the map $T_1(X)=\tr(X)\Id-X$ on $\MN{D}$ is positive.
\end{theorem}

\begin{proof}\leanok
    If $X\ge 0$, diagonalize $X$.  Its eigenvalues
    $\lambda_1,\ldots,\lambda_D$ are nonnegative, and for each $j$,
    \[
        \lambda_j\le \sum_{i=1}^D \lambda_i=\tr(X).
    \]
    Therefore all eigenvalues of $\tr(X)\Id-X$ are nonnegative.
\end{proof}

\begin{theorem}[Self-duality of the reduction map]
    \label{thm:reduction_map_self_dual}
    \lean{Matrix.traceAdjointMap_reductionMap}
    \leanok
    \uses{def:reduction_map}
    The reduction map is self-dual for the trace pairing:
    \[
        \tr(T_k(\rho)X)=\tr(\rho T_k(X)).
    \]
\end{theorem}

\begin{proof}\leanok
    Expand both sides using
    $T_k(Y)=\tr(Y)\Id-k^{-1}Y$ and bilinearity of the trace.  Both sides are
    \[
        \tr(\rho)\tr(X)-k^{-1}\tr(\rho X).
    \]
\end{proof}

\begin{theorem}[Choi matrix of the reduction map]
    \label{thm:choi_matrix_reduction_map}
    \lean{ChoiJamiolkowski.choiMatrix_reductionMap}
    \leanok
    \uses{def:reduction_map}
    For $D\ge 1$, the Choi matrix of
    $T_k(X)=\tr(X)\Id-k^{-1}X$ is
    \[
        \tau(T_k)=D^{-1}\Id-k^{-1}\ket{\Omega}\bra{\Omega}.
    \]
\end{theorem}

\begin{proof}\leanok
    The trace term sends the $(a,b)$ slice of
    $\ket{\Omega}\bra{\Omega}$ to its trace times the identity, giving
    the contribution $D^{-1}\Id$.  The second term is the Choi matrix of
    $k^{-1}$ times the identity map, namely
    $k^{-1}\ket{\Omega}\bra{\Omega}$.
\end{proof}

\begin{theorem}[Conjugation filters preserve positivity]
    \label{thm:conjugation_filters_preserve_positivity}
    \lean{unitaryConjLM_isCPMap, unitaryConjLM_isPositiveMap,
        IsPositiveMap.comp_unitaryConjLM, unitaryConjLM_isTP_of_unitary,
        unitaryConjLM_isChannel_of_unitary}
    \leanok
    \uses{def:cp_map, def:positive_map, def:trace_preserving, def:channel}
    For every $X\in\MN{D}$, the conjugation filter
    \[
        \rho\longmapsto X\rho X^\dagger
    \]
    is completely positive, hence positive.  Therefore, if
    $T:\MN{D}\to\MN{D}$ is positive, then
    $\rho\mapsto T(X\rho X^\dagger)$ is positive.  If in addition
    $X^\dagger X=\Id$, then the conjugation filter is trace-preserving and
    hence is a quantum channel.
\end{theorem}

\begin{proof}\leanok
    Complete positivity follows from the single Kraus operator $X$.
    Positivity follows immediately.  If $X^\dagger X=\Id$, cyclicity of trace
    gives
    \[
        \tr(X\rho X^\dagger)=\tr(X^\dagger X\rho)=\tr(\rho).
    \]
\end{proof}

\begin{definition}[Positive map between matrix algebras]
    \label{def:positive_map_rectangular}
    A linear map $T : \MN{D} \to \MN{D'}$ between matrix algebras of
    possibly different dimensions is \emph{positive} if $T(X) \ge 0$
    whenever $X \ge 0$, where $X \ge 0$ means that $X$ is positive
    semidefinite.  For $D' = D$ this is the notion of
    Definition~\ref{def:positive_map}.
\end{definition}

\begin{definition}[Trace-preserving map between matrix algebras]
    \label{def:trace_preserving_rectangular}
    A linear map $T : \MN{D} \to \MN{D'}$ between matrix algebras of
    possibly different dimensions is \emph{trace-preserving} if
    $\tr(T(X)) = \tr(X)$ for all $X$.  For $D' = D$ this is the notion of
    Definition~\ref{def:trace_preserving}.
\end{definition}

\begin{lemma}[Trace-normalization criterion, trace-preservation component]
    \label{lem:trace_normalization_criterion_trace}
    \lean{trace_comp_unitaryConjLM_of_conj_traceAdjointMap_one}
    \leanok
    \uses{def:trace_preserving_rectangular, def:trace_pairing_adjoint}
    Let $T:\MN{D}\to\MN{D'}$ be a linear map and let $X\in\MN{D}$ satisfy
    $X^\dagger T^*(\Id)X=\Id$.  Then for every $\rho\in\MN{D}$,
    \[
        \tr\!\bigl(T(X\rho X^\dagger)\bigr) = \tr(\rho).
    \]
\end{lemma}

\begin{proof}\leanok
    Expanding with the trace-pairing adjoint,
    $\tr(T(X\rho X^\dagger))=\tr(X\rho X^\dagger\,T^*(\Id))
    =\tr(\rho\,X^\dagger T^*(\Id)X)=\tr(\rho)$.
\end{proof}

\begin{theorem}[Filtered trace-normalization criterion]
    \label{thm:filtered_trace_normalization_criterion}
    \lean{comp_unitaryConjLM_positive_tracePreserving_of_conj_traceAdjointMap_one}
    \leanok
    \uses{def:positive_map_rectangular, def:trace_preserving_rectangular,
        def:trace_pairing_adjoint, lem:trace_normalization_criterion_trace,
        thm:conjugation_filters_preserve_positivity}
    Let $T:\MN{D}\to\MN{D'}$ be a positive map between matrix algebras of
    possibly different dimensions, and let $X\in\MN{D}$ satisfy
    \[
        X^\dagger T^*(\Id)X=\Id,
    \]
    where $T^*$ is the trace-pairing adjoint.  Then
    \[
        \rho\longmapsto T(X\rho X^\dagger)
    \]
    is positive and trace-preserving.

    This is the algebraic trace-normalization step in the proof of
    \cite[Chapter~3, Lemma ``Making positive maps trace preserving'']
    {Wolf2012Quantum}.  The inverse-square-root choice is recorded in the
    next theorem.
\end{theorem}

\begin{proof}\leanok
    Positivity follows because $X\rho X^\dagger$ is positive whenever
    $\rho$ is positive, and $T$ is positive.  For trace preservation, use the
    trace-pairing adjoint identity and cyclicity of trace:
    \[
        \tr(T(X\rho X^\dagger))
        =\tr(T^*(\Id)X\rho X^\dagger)
        =\tr(X^\dagger T^*(\Id)X\rho)
        =\tr(\rho).
    \]
\end{proof}

\begin{lemma}[Normalizing conjugation by inverse square root]
    \label{lem:conj_inv_sqrt_normalizes}
    \lean{Matrix.conjTranspose_inv_sqrt_mul_self_mul_inv_sqrt_eq_one_of_posDef}
    \leanok
    \uses{}
    Let $A\in\MN{D}$ be positive definite and set $S=A^{1/2}$.  Then
    \[
        (S^{-1})^\dagger A S^{-1}=\Id .
    \]
\end{lemma}

\begin{proof}\leanok
    Since $A$ is positive definite, $S$ is invertible and self-adjoint.
    Moreover $S^2=A$.  Therefore
    \[
        (S^{-1})^\dagger A S^{-1}
        =
        S^{-1}S^2S^{-1}
        =
        \Id .
    \]
\end{proof}

\begin{lemma}[Inverse square root of a positive definite matrix is invertible]
    \label{lem:pos_def_inv_sqrt_invertible}
    \lean{Matrix.PosDef.isUnit_det_inv_sqrt}
    \leanok
    \uses{}
    If $A\in\MN{D}$ is positive definite, then $(A^{1/2})^{-1}$ is
    invertible.
\end{lemma}

\begin{proof}\leanok
    A positive definite matrix is invertible, hence so is its square root
    $A^{1/2}$, and the inverse of an invertible matrix is invertible.
\end{proof}

\begin{theorem}[Trace normalization by inverse square root]
    \label{thm:trace_normalization_inverse_square_root}
    \lean{comp_unitaryConjLM_inv_cfc_sqrt_traceAdjointMap_one}
    \leanok
    \uses{def:positive_map_rectangular, def:trace_preserving_rectangular,
        def:trace_pairing_adjoint,
        thm:filtered_trace_normalization_criterion, lem:conj_inv_sqrt_normalizes}
    Let $T:\MN{D}\to\MN{D'}$ be a positive map such that $T^*(\Id)$ is positive
    definite.  Put
    \[
        X=(T^*(\Id))^{-1/2}.
    \]
    Then
    \[
        \rho\longmapsto T(X\rho X^\dagger)
    \]
    is positive and trace-preserving.
\end{theorem}

\begin{proof}\leanok
    Let $S=(T^*(\Id))^{1/2}$.  Since $T^*(\Id)$ is positive definite, $S$ is
    invertible and self-adjoint, and $S^2=T^*(\Id)$.  For $X=S^{-1}$,
    \[
        X^\dagger T^*(\Id)X
        =
        S^{-1}S^2S^{-1}
        =
        \Id .
    \]
    The preceding trace-normalization criterion applies.
\end{proof}

\begin{lemma}[Making positive maps trace preserving]
    \label{lem:making_positive_maps_trace_preserving}
    \lean{exists_isUnit_det_comp_unitaryConjLM_positive_tracePreserving}
    \leanok
    \uses{def:positive_map_rectangular, def:trace_preserving_rectangular,
        def:trace_pairing_adjoint,
        thm:trace_normalization_inverse_square_root,
        lem:pos_def_inv_sqrt_invertible}
    Let $T:\MN{D}\to\MN{D'}$ be a positive map for which $T^*(\Id)$ is
    positive definite.  Then there exists an invertible $X\in\MN{D}$ such
    that
    \[
        \rho\longmapsto T(X\rho X^\dagger)
    \]
    is a trace-preserving positive map
    \cite[Chapter~3, Lemma ``Making positive maps trace preserving'']
    {Wolf2012Quantum}.
\end{lemma}

\begin{proof}\leanok
    Take $X=(T^*(\Id))^{-1/2}$, which is invertible by
    Lemma~\ref{lem:pos_def_inv_sqrt_invertible} because $T^*(\Id)$ is
    positive definite.  By the preceding theorem,
    $\rho\mapsto T(X\rho X^\dagger)$ is positive and trace-preserving.
\end{proof}

\begin{theorem}[The trace bounds a positive matrix]
    \label{thm:psd_le_trace_smul_one}
    \lean{Matrix.PosSemidef.trace_smul_one_sub_self_posSemidef}
    \leanok
    Let $D\geq 1$ and let $P\in\MN{D}$ be positive semidefinite. Then
    \[
        P\preceq \tr(P)\Id.
    \]
\end{theorem}

\begin{proof}\leanok
    Diagonalize $P$ with nonnegative eigenvalues $\lambda_1,\ldots,\lambda_D$.
    Each eigenvalue is bounded above by their sum $\tr(P)$, so every eigenvalue
    of $\tr(P)\Id-P$ is nonnegative.
\end{proof}

\begin{theorem}[Forward Lorentz-cone trace inequality]
    \label{thm:psd_trace_square_le_square_trace}
    \lean{Matrix.PosSemidef.trace_sq_re_le_trace_re_sq}
    \leanok
    If $A \in \MN{D}$ is positive semidefinite, then
    \[
        \re\tr(A^2) \le \left(\re\tr(A)\right)^2 .
    \]
\end{theorem}

\begin{proof}\leanok
    Diagonalize $A$ with nonnegative eigenvalues $\lambda_i$. Then
    $\tr(A^2)=\sum_i \lambda_i^2$ and $\tr(A)=\sum_i \lambda_i$, so the
    desired inequality is
    \[
        \sum_i \lambda_i^2 \le \left(\sum_i \lambda_i\right)^2 ,
    \]
    which follows because all mixed products are nonnegative.
\end{proof}

\begin{theorem}[Trace-nonnegative Lorentz-cone converse]
    \label{thm:trace_nonnegative_lorentz_cone_converse}
    \lean{Matrix.IsHermitian.posSemidef_of_trace_re_nonneg_of_card_sub_one_mul_trace_sq_re_le}
    \leanok
    Let $A\in\MN{D}$ be Hermitian and suppose that
    $\re\tr(A)\ge0$. If
    \[
        (D-1)\re\tr(A^2)
        \le \left(\re\tr(A)\right)^2 ,
    \]
    then $A$ is positive semidefinite.

    This is the trace-nonnegative form of the converse in
    \cite[Proposition~3.9]{Wolf2012Quantum}. The unrestricted squared
    implication printed there is false without a trace sign condition; this is
    recorded in
    \path{docs/paper-gaps/wolf_ch3_lorentz_cone_trace_sign.tex}.
\end{theorem}

\begin{proof}\leanok
    Diagonalize $A$ with real eigenvalues $\lambda_i$. Suppose that some
    eigenvalue $\lambda_0$ is negative. Let
    \[
        T=\sum_i\lambda_i,\qquad R=\sum_{i\ne 0}\lambda_i .
    \]
    Since $T\ge0$ and $T=\lambda_0+R$, one has $T<R$, and hence $T^2<R^2$.
    Cauchy's inequality gives
    \[
        R^2\le (D-1)\sum_{i\ne 0}\lambda_i^2
            \le (D-1)\sum_i\lambda_i^2,
    \]
    contradicting the assumed Lorentz inequality. Thus all eigenvalues are
    nonnegative.
\end{proof}

\begin{definition}[Sum of the largest eigenvalues]
    \label{def:ky_fan_norm}
    \lean{Matrix.IsHermitian.kyFanNorm}
    \leanok
    Let $A\in\MN{D}$ be Hermitian with eigenvalues
    $\lambda_1\ge\lambda_2\ge\cdots\ge\lambda_D$ listed in decreasing order. For
    $0\le k\le D$ set
    \[
        S_k(A) = \sum_{i=1}^{k}\lambda_i ,
    \]
    the sum of the $k$ largest (signed) eigenvalues. Indices beyond $D$ contribute
    zero, so the value stabilizes at $\re\tr(A)$ once $k$ reaches
    $D$.

    The \emph{Ky-Fan $k$-norm} $\|A\|_{(k)}$ is by definition the sum of the $k$
    largest \emph{singular values} of $A$. It agrees with $S_k(A)$ exactly when
    $A$ is positive semidefinite, where eigenvalues and singular values coincide;
    for an indefinite Hermitian $A$ the two differ (for $A=\diag(1,-2)$ one has
    $S_1(A)=1$ but $\|A\|_{(1)}=2$). The results below concern $S_k(A)$ for
    arbitrary Hermitian $A$, so on the positive-semidefinite cone they are
    statements about the Ky-Fan norm.
\end{definition}

\begin{theorem}[Achievability of the largest-eigenvalue sum]
    \label{thm:ky_fan_achievability}
    \lean{Matrix.IsHermitian.exists_isProj_trace_eq_kyFanNorm}
    \leanok
    \uses{def:ky_fan_norm}
    For every Hermitian $A\in\MN{D}$ and every $k\ge0$ there is an orthogonal
    projection $P$ of rank $\min(k,D)$ with
    \[
        \re\tr(PA) = S_k(A) .
    \]
\end{theorem}

\begin{proof}\leanok
    \uses{def:ky_fan_norm}
    Diagonalize $A=U\,\diag(\lambda_i)\,U^\dagger$ and take $P$ to be the
    projection onto the span of the eigenvectors with the $\min(k,D)$ largest
    eigenvalues. Then $\re\tr(PA)=\sum_{i\le k}\lambda_i=S_k(A)$.
\end{proof}

\begin{theorem}[Upper bound for the largest-eigenvalue sum]
    \label{thm:ky_fan_upper_bound}
    \lean{Matrix.IsHermitian.trace_mul_re_le_kyFanNorm}
    \leanok
    \uses{def:ky_fan_norm}
    For every Hermitian $A\in\MN{D}$ and every orthogonal projection $P$ of rank
    $k<D$,
    \[
        \re\tr(PA) \le S_k(A) .
    \]
\end{theorem}

\begin{proof}\leanok
    \uses{def:ky_fan_norm}
    Conjugating $P$ by $U$ gives an orthogonal projection $Q$ of the same rank,
    and $\re\tr(PA)=\sum_i w_i\lambda_i$ where $w_i$ are the
    diagonal entries of $Q$. Each $w_i$ lies in $[0,1]$ and the weights sum to
    $\tr P = k$. With the threshold $c=\lambda_{k+1}$,
    \[
        \sum_{i\le k}\lambda_i - \sum_i w_i\lambda_i
            = \sum_{i\le k}(\lambda_i-c)(1-w_i)
                + \sum_{i>k}(c-\lambda_i)w_i \ge 0 ,
    \]
    since for $i\le k$ the eigenvalues dominate $c$ while $1-w_i\ge0$, and for
    $i>k$ they are dominated by $c$ while $w_i\ge0$.
\end{proof}

\begin{theorem}[Ky Fan's maximum principle]
    \label{thm:ky_fan_maximum_principle}
    \lean{Matrix.IsHermitian.kyFanNorm_eq_sup_trace}
    \leanok
    \uses{def:ky_fan_norm, thm:ky_fan_achievability, thm:ky_fan_upper_bound}
    Let $A\in\MN{D}$ be Hermitian and $0\le k<D$. Then the sum of its $k$ largest
    eigenvalues is the largest value of $\re\tr(PA)$ attained by an
    orthogonal projection $P$ of rank $k$:
    \[
        S_k(A) = \max_{P^2=P=P^\dagger,\ \tr P = k}
            \re\tr(PA) .
    \]
    This is Ky Fan's maximum principle \cite{Fan1949Theorem}; see also
    \cite[Problem~I.6.15, Exercise~II.1.13]{Bhatia1997Matrix}. For a positive
    semidefinite $A$, where $S_k(A)$ coincides with the Ky-Fan norm
    $\|A\|_{(k)}$, it gives the extremal overlap of \cite[Lemma~3.1]{Wolf2012Quantum}.
\end{theorem}

\begin{proof}\leanok
    \uses{thm:ky_fan_achievability, thm:ky_fan_upper_bound}
    The achievability half (Theorem~\ref{thm:ky_fan_achievability}) provides a
    rank-$k$ projection whose trace against $A$ equals $S_k(A)$, so the value is
    attained. The upper bound (Theorem~\ref{thm:ky_fan_upper_bound}) shows that no
    rank-$k$ projection exceeds it. Hence $S_k(A)$ is the maximum.
\end{proof}

\begin{definition}[Maps preserving the positive semidefinite cone]
    \label{def:maps_psd_cone_onto}
    \lean{MapsPSDConeOnto}
    \leanok
    \uses{def:positive_map}
    A linear map $T:\MN{D}\to\MN{D}$ \emph{maps the positive semidefinite cone
    onto itself} if $T$ is positive and every positive semidefinite matrix is
    the image under $T$ of a positive semidefinite matrix. This is condition
    (1) of \cite[Proposition~3.8]{Wolf2012Quantum}.
\end{definition}

\begin{theorem}[Conjugations preserve the positive semidefinite cone]
    \label{thm:conjugation_maps_psd_cone_onto}
    \lean{unitaryConjLM_mapsPSDConeOnto,
        unitaryConjLM_comp_transpose_mapsPSDConeOnto}
    \leanok
    \uses{def:maps_psd_cone_onto}
    Let $Y\in\MN{D}$ be invertible. Then the map $X\mapsto YXY^\dagger$ and the
    map $X\mapsto YX^TY^\dagger$ each map the positive semidefinite cone onto
    itself. This is the implication (3) $\Rightarrow$ (1) of
    \cite[Proposition~3.8]{Wolf2012Quantum}.
\end{theorem}

\begin{proof}\leanok
    Conjugation by any matrix preserves positive semidefiniteness, and the
    transpose of a positive semidefinite matrix is positive semidefinite, so
    both maps send the cone into itself. For surjectivity onto the cone, given
    a positive semidefinite $A$ set $W=Y^{-1}A(Y^{-1})^\dagger$, which is again
    positive semidefinite. Then $YWY^\dagger=A$ proves the conjugation case,
    and $Y(W^T)^TY^\dagger=A$ proves the transpose case, with $W$ and $W^T$
    positive semidefinite.
\end{proof}

\begin{theorem}[Transposition is not completely positive]
    \label{thm:transpose_not_completely_positive}
    \lean{ChoiJamiolkowski.transposeLinearMapComplex_not_isCPMap}
    \leanok
    \uses{def:cp_map, def:flip_operator, def:fin_two_antisymm_vec,
        thm:choi_transpose_flip, thm:cp_iff_choi_posSemidef}
    If $D \ge 2$, the matrix transposition map $\theta(X)=X^T$ on $\MN{D}$
    is not completely positive.
\end{theorem}

\begin{proof}\leanok
    \uses{def:flip_operator, def:fin_two_antisymm_vec, thm:choi_transpose_flip,
        thm:cp_iff_choi_posSemidef}
    If $\theta$ were completely positive, its Choi matrix would be positive
    semidefinite. By Theorem~\ref{thm:choi_transpose_flip}, this Choi matrix is
    $D^{-1}F$. For the antisymmetric vector
    $v=\ket{01}-\ket{10}$,
    \[
        \bra{v}\,D^{-1}F\,\ket{v} = -\frac{2}{D} < 0,
    \]
    contradicting positive semidefiniteness.
\end{proof}

\begin{theorem}[Entrywise complete positivity from a Kraus representation]
    \label{thm:cp_map_entrywise_complete_positive}
    \lean{IsCPMap.map_cstarMatrix_nonneg}
    \leanok
    \uses{def:cp_map}
    Let $E : \MN{D} \to \MN{D}$ be completely positive in the Kraus sense.
    For every $k$ and every positive block matrix
    $M \in M_k(\MN{D})$, the entrywise image
    \[
        (E(M_{ab}))_{a,b}
    \]
    is again positive in $M_k(\MN{D})$.
\end{theorem}

\begin{proof}\leanok
    Write $E(X)=\sum_i K_i X K_i^\dagger$.  For each $i$, let $d_i$ be
    the block-diagonal matrix with $K_i$ on every diagonal block.  Then
    \[
        (E(M_{ab}))_{a,b}=\sum_i d_i M d_i^\dagger .
    \]
    Each term on the right is positive because conjugation preserves
    positivity, and a finite sum of positive matrices is positive.
\end{proof}

\begin{theorem}[Kraus maps as abstract completely positive maps]
    \label{thm:cp_map_to_abstract_completely_positive}
    \lean{IsCPMap.toCompletelyPositiveMap}
    \leanok
    \uses{def:cp_map, thm:cp_map_entrywise_complete_positive}
    Every Kraus-represented completely positive map
    $E : \MN{D} \to \MN{D}$ determines a completely positive map between
    the matrix $C^*$-algebras in the entrywise positivity sense.
\end{theorem}

\begin{proof}\leanok
    The defining entrywise positivity condition is exactly
    Theorem~\ref{thm:cp_map_entrywise_complete_positive}.
\end{proof}

\begin{theorem}[CP maps are positive]\label{thm:cp_is_positive}
    \lean{IsCPMap.isPositiveMap}
    \leanok
    \uses{def:cp_map, def:positive_map}
    Every completely positive map is positive.
\end{theorem}

\begin{proof}\leanok
    If $X \ge 0$, then each summand $K_i X K_i^\dagger \ge 0$
    (conjugation preserves positive semidefiniteness),
    so the sum is PSD.
\end{proof}

\begin{theorem}[Channels are positive]\label{thm:channel_positive}
    \lean{IsChannel.pos}
    \leanok
    \uses{def:channel, def:positive_map}
    Every quantum channel is positive.
\end{theorem}

\begin{proof}\leanok\uses{thm:cp_is_positive}
    A quantum channel is completely positive by definition, so
    Theorem~\ref{thm:cp_is_positive} applies.
\end{proof}

\begin{theorem}[Positive maps preserve Hermiticity]
    \label{thm:positive_map_preserves_hermitian}
    \lean{IsPositiveMap.map_isHermitian}
    \leanok
    \uses{def:positive_map_is_positive_linear_map}
    If $E:\MN{n}\to\MN{m}$ is positive and
    $X\in\MN{n}$ satisfies $X = X^\dagger$, then
    $E(X) = E(X)^\dagger$ in $\MN{m}$.
\end{theorem}

\begin{proof}\leanok
    Regard $E$ as the associated positive linear map
    (Definition~\ref{def:positive_map_is_positive_linear_map}). For positive
    linear maps between matrix $C^*$-algebras, $E(Y^\ast)=E(Y)^\ast$.
    Since $X=X^\dagger$ means $X=X^\ast$, one obtains
    $E(X)^\dagger=E(X)^\ast=E(X^\ast)=E(X)$.
\end{proof}

\begin{definition}[Quantum channel]\label{def:channel}
    \lean{IsChannel}
    \leanok
    \uses{def:cp_map, def:trace_preserving}
    A linear map $E : \MN{D} \to \MN{D}$ is a \emph{quantum channel}
    if it is both completely positive and trace-preserving (CPTP),
    following~\cite{Wolf2012Quantum}.
\end{definition}

\begin{definition}[Density matrices]\label{def:density_matrices}
    \lean{densityMatrices, mem_densityMatrices}
    \leanok
    The set of \emph{density matrices} in $\MN{D}$ is
    \[
        \mc{D}_D =  \{\rho \in \MN{D} : \rho \ge 0,  \tr(\rho) = 1\}.
    \]
\end{definition}

\begin{theorem}[Density matrices are compact]\label{thm:density_compact}
    \lean{densityMatrices_isCompact}
    \leanok
    \uses{def:density_matrices}
    $\mc{D}_D$ is compact (in the entry-wise topology on $\MN{D}$).
\end{theorem}

\begin{proof}
    \leanok
    The PSD cone is closed
    (preimage of the closed non-negative cone under continuous
    quadratic forms), and if $\rho \ge 0$ with $\tr(\rho) = 1$
    then each entry satisfies $|\rho_{ij}|^2 \le \rho_{ii}\rho_{jj} \le 1$:
    the diagonal entries are non-negative and sum to~$1$, and the
    off-diagonal bound is the PSD Cauchy--Schwarz inequality.
    Hence the matrix entries are uniformly bounded, and Heine--Borel
    gives compactness.
\end{proof}

\begin{theorem}[Density matrices are convex]\label{thm:density_convex}
    \lean{densityMatrices_isConvex}
    \leanok
    \uses{def:density_matrices}
    $\mc{D}_D$ is convex.
\end{theorem}

\begin{proof}\leanok
    Convex combination of PSD matrices is PSD, and trace is linear.
\end{proof}

\begin{theorem}[Density matrices are nonempty]\label{thm:density_nonempty}
    \lean{densityMatrices_nonempty}
    \leanok
    \uses{def:density_matrices}
    For $D \ge 1$, $\mc{D}_D \neq \varnothing$.
\end{theorem}

\begin{proof}\leanok
    $\frac{1}{D} \Id_D$ is a density matrix.
\end{proof}

\begin{theorem}[Channels preserve density matrices]\label{thm:channel_preserves_density}
    \lean{IsChannel.map_densityMatrices}
    \leanok
    \uses{def:channel, def:density_matrices}
    If $E$ is a quantum channel, then $E(\mc{D}_D) \subseteq \mc{D}_D$.
\end{theorem}

\begin{proof}\leanok\uses{thm:cp_is_positive}
    Complete positivity implies positivity
    (Theorem~\ref{thm:cp_is_positive}),
    so $E(\rho) \ge 0$; trace preservation gives
    $\tr(E(\rho)) = \tr(\rho) = 1$.
\end{proof}
